% TeX_root = ../main.tex

\marginnote{\scriptsize 13/02/2025 }

\begin{theorem}[Bicommutant Theorem]
  Let $A \subset B(\mathcal{H})$ be a unital $*$-subalgebra, where
  $\mathcal{H}$ is separable. Then $\overline{ \mathcal{A}}^{SOT} =
  \mathcal{A}^{\prime \prime}$
\end{theorem}
\begin{proof}
  Let $T \in \overline{\mathcal{A}}^{SOT}$, and $S \in
  \mathcal{A}^\prime$. Let $(T_i) \in \mathcal{A}$ be a net
  converging in SOT to $T$. Then for all $\xi \in \mathcal{H}$,
  \begin{align*}
    TS \xi = \lim_{i} T_iS \xi = \lim_{i} ST_i \xi = S \lim_i T_i \xi = ST \xi
  \end{align*}
  shows that $T \in \mathcal{A}^{\prime\prime}$.

  To see the converse, let $T \in \mathcal{A}^{\prime \prime}$, fix
  $\xi \in \mathcal{H}$, and $\varepsilon >
  0$ and let $K = \overline{\mathcal{A} \xi}$. Then notice that $K$
  is a reducing subspace for all operators in $\mathcal{A}$.
  \textcolor{red}{verify}. Then $P_K T = TP_K$ by our results on
  reducing subsapces in last semester. Thus $P_K \in \mathcal{A}^\prime$.

  %Prof version
  Let $\xi_1 , \xi_2 , \ldots , \xi_n \in \mathcal{H}$, $\varepsilon>
  0$ be given. Let
  \begin{align*}
    \mathcal{K} = \Big \{
      \begin{bmatrix}%{c}
        A \xi_1 \\
        \vdots \\
        A \xi_n
    \end{bmatrix}  \ : \  A \in \mathcal{A} \Big \}
  \end{align*}
  For each $A \in \mathcal{A}$, let $\mathcal{A}^{(n)} = I_n \otimes A\in
  B(\mathcal{H}^n)$. Let $\mathcal{A}^n = \{ A^{(n)}  \ : \  A \in
  \mathcal{A} \}$. Observe that $\mathcal{A}^{(n)}(\mathcal{K})
  \subset \mathcal{K}$. Also observe that $\mathcal{K}$ is reducing
  for all $A^{(n)} \in \mathcal{A}^{(n)}$. So $P_{\mathcal{K}} \in
  (\mathcal{A}^{(n)})^\prime = M_n(\mathcal{A}^\prime)$.
  \textcolor{red}{verify the rest}.
\end{proof}

\begin{definition}
  For every $\xi, \eta \in \mathcal{H}$, denote $\omega_{\xi, \eta}
  \in B(\mathcal{H})^*$ such that
  \begin{align*}
    \omega_{\xi, \eta}(T) = \langle T \xi ,  \eta \rangle
  \end{align*}
\end{definition}

\begin{theorem}
  $B(\mathcal{H}) = \overline{\textrm{span}}\{ \omega_{\xi, \eta}
  \ : \    \xi, \eta \in \mathcal{H} \}^*$
\end{theorem}

\begin{proposition}
  For $T \in B(\mathcal{H})$, the following are equivalent.
  \begin{enumerate}[label=(\arabic*)]
    \item $\langle T \xi ,  \xi \rangle \ge 0, \forall \xi
      \in \mathcal{H}$
    \item $\exists S \in B(\mathcal{H})$ such that  $T = S^*S$
  \end{enumerate}
\end{proposition}
\begin{proof}
  $2 \implies 1$ is obvious.

  Now for the other one, it is clear that $T$ is self adjoint. So we
  get that $\Xi: C(\textrm{sp}(T)) \to B(\mathcal{H})$. We claim that
  $\textrm{sp}(T) \subset \mathbb{R}^{\ge 0}$, which then completes
  the proof by the identification of $T$ with the multiplication
  operator corresponding to the identity function in $C(\textrm{sp}(T))$.

  For the sake of contradiction, assume that $Y = \textrm{sp}(T) \cap
  \mathbb{R}^{< 0} \neq \emptyset$. Then let $P = \Xi(\chi_Y)$.
\end{proof}
